\section{Related work}
\label{sec:related}

\paragraph{Agent benchmarks for experimentation.}
BoxingGym gives language agents ten sequential experiments in generative probabilistic
environments and evaluates experimental design by expected information gain (EIG), followed
by prediction or model communication \citep{boxinggym}. Its reported learning curves are
especially relevant here: prediction usually improves with more observations, but error
increases in the hyperbolic-discounting environment. \slowlab{} asks a different question.
The endpoint is an operating recommendation, the design is an allocation to named physical
units, and trials consume a fixed campaign budget and opportunity cost. We therefore score
decision regret and condition design diagnostics on the same history used for the
recommendation, rather than equating parameter information with decision value.

CARE also evaluates an LLM inside a budgeted sequential reaction-optimization loop
\citep{care2026}. It has the model write an executable ranking program while a numerical
reference policy and intervention gate retain control over experiment selection. Its matched
offline replay reveals one condition at a time from an enumerated pool. \slowlab{} instead
evaluates direct agent control of feasible batches, including allocation, correlated
measurements and a recommendation that need not be one of the tested treatments. These are
complementary tests of how an LLM enters an experimental controller.

\paragraph{Executable experiment-optimization benchmarks.}
Summit provides mechanistic and data-driven reaction benchmarks and common optimizer
interfaces \citep{summit}; Olympus adds noisy emulators trained on experimental chemistry
and materials datasets and a broad planner library \citep{olympus}. Minerva goes further
toward laboratory execution, supporting noisy multi-objective optimization in large
parallel batches and enforcing batch constraints in automated high-throughput chemistry
\citep{minerva}. These systems establish that noise, finite data and parallel execution are
already central benchmark concerns. \slowlab{} contributes a different resource geometry:
controls act at chamber and loop levels, so the agent must assign treatments to units under
split-plot constraints, and every tested treatment is charged by its economic shortfall.
Table~\ref{tab:related-comparison} uses ``not reported'' where the cited evaluation does not
specify a property; absence of a report is not evidence that a system cannot represent it.

\begin{table*}[t]
\centering\small
\setlength{\tabcolsep}{4pt}
\begin{tabularx}{\textwidth}{@{}l*{5}{>{\centering\arraybackslash}X}@{}}
\toprule
System & noisy outcomes & finite campaign & parallel / batch constraint & time-stamped interim data & separate final recommendation \\
\midrule
BoxingGym & yes & yes (10 queries) & not reported & not reported & yes (prediction/model) \\
Summit & not reported & user-set & not reported & not reported & N/A \\
Olympus & yes & user-set & not reported & not reported & N/A \\
Minerva & yes & yes & yes & not reported & N/A \\
CARE & dataset outcomes & yes (15 total) & N/A (sequential replay) & N/A & best tested condition \\
\slowlab{} & yes & yes & tiered physical capacity & yes & yes (continuous action) \\
\bottomrule
\end{tabularx}
\caption{Scope of the evaluated interfaces in the cited work. ``N/A'' means the published
task has a different endpoint; ``not reported'' avoids inferring a negative from an omitted
feature.}
\label{tab:related-comparison}
\end{table*}

\paragraph{Budget, delay and experimental structure.}
Budgeted Bayesian optimization studies how to allocate a total evaluation-cost budget,
including the non-myopic case where costs vary and are initially unknown
\citep{astudillo2021budgeted}. Delayed-feedback BO studies how to choose new evaluations
while earlier ones remain pending \citep{verma2022delayed}. These methods provide algorithmic
baselines and regret analyses; our benchmark fixes a domain-specific facility and exposes
the resulting control hierarchy to general-purpose agents. The chamber/loop hierarchy is a
split-plot design: whole-plot factors have fewer independent experimental units than
subplot factors, so analysis and randomization must respect two error strata
\citep{goos2011splitplot}. This statistical structure motivates both the feasibility checker
and the block-aware readers.

\paragraph{Decision-theoretic design.}
Bayesian experimental design commonly scores parameter information
\citep{lindley1956,chaloner1995}; decision-theoretic design instead values how observations
change a downstream action \citep{raiffa1961,howard1966}. Related quantities appear as the
knowledge gradient and entropy-search criteria in Bayesian optimization
\citep{frazier2008kg,hennig2012,hernandez2014}. Our common-history diagnostic applies this
decision view to benchmark traces: it separates risk supported by the observed history from
excess risk introduced by the submitted recommendation, while reporting the numerical
calibration limits of the posterior evaluator.

\paragraph{Scientific agents and self-driving laboratories.}
Scientific-agent suites often start from existing data, while autonomous chemistry and
materials systems connect planners to laboratory tools
\citep{scienceagentbench,discoverybench,coscientist,chemcrow,
szymanski2023autonomous,burger2020mobile}. CropGym-style environments expose crop simulators
as control problems \citep{wofostgym}. \slowlab{} occupies a narrower intersection: it is a
simulated benchmark for experimental allocation and final recommendation in one greenhouse
domain. We do not infer cross-domain validity from that instantiation.
